% Section 3: Bag-of-words representation and notation
% Numbers are from the screened sample (derived_screened/, globals.json BOW_* settings), 2026-11-06.

\section{Methods}

\subsection{Notation}

Let $\mathcal{D} = \{d_1, \dots, d_N\}$ denote the corpus, where each document $d_n$ is the
cleaned auto-caption transcript of one school-board meeting video. In our screened 2018--2023
sample $N = 80{,}063$. Each document is a sequence of tokens
\[
  d_n = (w_{n1}, w_{n2}, \dots, w_{n L_n}), \qquad w_{n\ell} \in \mathcal{V},
\]
where $L_n$ is the length of document $n$ and $\mathcal{V} = \{v_1, \dots, v_V\}$ is the
vocabulary of $V = |\mathcal{V}|$ distinct terms. Each document belongs to a district
$i(n) \in \{1, \dots, I\}$ ($I = 1{,}485$ districts, indexed by NCES LEAID) and a time period
$t(n)$ (calendar year or quarter of the meeting date). Each district has a partisan label
$P_i \in \{\mathsf{R}, \mathsf{D}\}$, equal to $\mathsf{R}$ if its county's 2020 two-party
Republican presidential vote share exceeds one half. To avoid a clash with the corpus
$\mathcal{D}$, we write the parties in sans-serif and the sets of districts as
$\mathcal{I}^{\mathsf{R}}$ and $\mathcal{I}^{\mathsf{D}}$.

\subsection{Bag-of-words representation}

The bag-of-words model discards word order and keeps only how often each term occurs.
Document $d_n$ is represented by its count vector
\[
  \mathbf{c}_n = (c_{n1}, \dots, c_{nV}) \in \mathbb{N}^{V}, \qquad
  c_{nj} = \sum_{\ell=1}^{L_n} \mathbf{1}\{w_{n\ell} = v_j\},
\]
and the corpus by the $N \times V$ document--term matrix $\mathbf{C} = [c_{nj}]$. Row sums
$m_n = \sum_j c_{nj}$ give document length after preprocessing, and normalised frequencies are
$\mathbf{q}_n = \mathbf{c}_n / m_n$.

\paragraph{Preprocessing.} Captions are raw YouTube auto-transcriptions with no punctuation,
casing or speaker labels. We remove timestamps and bracketed tags, lowercase, and strip
punctuation (keeping in-word apostrophes). We then drop English stopwords, contractions,
caption filler (\emph{um}, \emph{uh}) and procedural meeting language (\emph{motion},
\emph{second}, \emph{aye}), and keep only terms appearing in at least 20 transcripts. This gives
$V = 85{,}890$ terms and about 399 million tokens. For the partisanship analysis we further drop
numerals, terms under three characters and terms used in fewer than 25 districts, and merge
possessives, leaving $V = 54{,}833$.

\paragraph{Aggregation.} The district, not the meeting, is our unit of analysis, because
partisanship is measured at the district (county) level and districts differ widely in how many
meetings they post. We sum counts within district $i$ and period $t$:
\[
  \mathbf{c}_{it} = \sum_{n:\, i(n) = i,\ t(n) = t} \mathbf{c}_n, \qquad
  m_{it} = \sum_j c_{itj}, \qquad
  \mathbf{q}_{it} = \mathbf{c}_{it} / m_{it},
\]
which yields 18{,}898 district--quarter units. Every later method (GST partisanship, TF-IDF, the
culture-war lexicon) is a function of $\{\mathbf{c}_{it}\}$, and each district receives equal
weight regardless of how many meetings it posted.

\paragraph{Why bag of words.} Our questions are about \emph{what} boards talk about and
\emph{which words} distinguish Republican from Democratic districts, not about syntax. Counts
are well suited to this: they are transparent, scale to $\sim\!10^5$ long documents, and are the
input assumed by both the Gentzkow--Shapiro--Taddy (2019) partisanship estimator and topic
models. They are also robust to the main weakness of our text: auto-captions garble punctuation
and sentence boundaries, but most content words are transcribed correctly. The cost is that a
count records that a word was said, not its stance or speaker; a mention of ``mask mandate''
is the same whether it is made in support or in protest. We return to this in the discussion.
